\section{Limitations and Threats to Validity}\label{sec:limits}
The evidence in this study is subject to several specific constraints that delimit its generalizability beyond the evaluated architecture.

The evaluation relies on deterministic replay of 64 fixed technical-interview instances organized into eight structural clusters and evaluated under 20 question-order seeds ($N = 160$ sessions per condition). While cluster-bootstrap resampling (Field and Welsh) accounts for this dependency, the empirical sample remains bounded by the eight clusters, and leave-one-cluster-out analysis shows that the magnitude of the level-change effect is sensitive to cluster composition. In particular, the observed increase of $+0.59$ proactive probe turns per session is concentrated entirely in replay cluster $k=1$; in the remaining seven clusters, proactive probing added exactly 0.00 turns. Replay dynamics cannot be assumed to generalize to broader interview populations.

Because candidate responses were pre-computed constants, the replay design omits dynamic candidate adaptation. Changing question difficulty or inserting a follow-up probe could not elicit an adapted response from a human participant. While this isolates coordination mechanics from human variance, it precludes evaluating closed-loop feedback effects on candidate learning, retention, confidence, or placement success. The trajectory metrics describe system decision dynamics under fixed evaluator inputs, not pedagogical efficacy.

The behavioral equivalence observed between specialized and unified architectures ($\mathbb{P}(\text{agree}) = 1.0$) applies strictly to deterministic rule- and ledger-based implementations under open-loop replay (1,280 turns). It does not establish equivalence for stochastic agents, live large language models, prompt variations, or closed-loop environments where agent actions alter subsequent observation streams. Furthermore, the Strategy and Coaching agents operate on explicit heuristic rules rather than autonomous learning, and the absence of live generative LLMs precludes assessing language-driven coordination or open-ended failure modes.

The adaptive baseline is represented by a single frozen PPO checkpoint (\texttt{seed\_123}); observed trajectory dynamics reflect this specific policy rather than general properties of reinforcement learning. The arbitration threshold $\tau = 0.60$ is a heuristic engineering parameter rather than a mathematically optimal value. Because replay payloads lacked raw acoustic waveforms, acoustic confidence remained at neutral defaults, leaving both the stretch-proposal branch and the synthetic grid regions (11.8\%) where deployed and canonical guardrails disagree unexercised.

Dataset provenance introduces additional boundaries. The 64 replay instances combine constructed and team-generated payloads. The replay experiments treat the deployed evaluator outputs as fixed inputs; evaluator validity is not independently evaluated in this study, and no human ratings are used in the reported experiments. Evaluator inaccuracies would therefore propagate into the state representation and downstream decisions.

No demographic fairness, adverse impact, or psychometric subgroup analyses were conducted, and no formal institutional ethics review was obtained; the offline replay corpus contains no personal data and is used strictly for software evaluation. Finally, the fault-injection experiments demonstrate containment only for the tested proposal defects and interfaces, without establishing universal fault tolerance or complete semantic correctness.
